\chapter{Introducción}

El presente documento tiene como objetivo presentar la planeación de la 
iteración~\textbf{\#\axSprintId} del proyecto \textbf{\axProjectName}, que se llevará a cabo 
entre el \textbf{\axSprintStart} y el \textbf{\axSprintEnd}. Este plan incluye la definición de 
los objetivos, las actividades a realizar, los recursos necesarios y los 
entregables esperados durante esta iteración.

\paragraph{}
Este documento es el resultado de aprobación de la iteración anterior y la 
retroalimentación recibida por parte del mandante, dueño del producto, líder de 
entrega, especialistas y terceras partes.

\section{Limitaciones}
Toda la promesa y el alcance de la iteración están sujetos a la disponibilidad de los
recursos y pueden verse afectados por factores imprevistos, incluyendo, pero no 
limitados a, eventos de fuerza mayor, interrupciones externas, condiciones 
climáticas extremas, problemas de infraestructura, decisiones regulatorias, o 
cualquier otro evento fuera de nuestro control razonable. Nos comprometemos a 
hacer todo lo posible para cumplir con los objetivos establecidos, pero no 
podemos garantizar su finalización si tales eventos ocurren.

\section{Nuestra filosofía}
Aspiramos a ofrecer a nuestros clientes dos cosas:

\begin{itemize}
  \item \textbf{Transparencia}: Que nuestro cliente sepa en todo momento lo que
  estamos haciendo, incluso cuando hay obstáculos en el camino.
  \item \textbf{Predictibilidad}: Que nuestro cliente tenga tranquilidad de que 
  estamos haciendo lo que dijimos que haríamos, y que si hay cambios, que tenemos 
  mecanismos para adaptarnos a ellos.
\end{itemize}

\subsection{¿Cómo organizamos nuestro trabajo?}
Nosotros descomponemos el trabajo de la siguiente forma:

\begin{itemize}
  \item \textbf{Programa}: Es un conjunto de proyectos relacionados entre sí. 
  Programas pequeños tienen un solo proyecto.
  \item \textbf{Proyecto}: Es un conjunto de hitos que se deben cumplir para 
  lograr un objetivo. En AX ofrecemos proyectos de arquitectura, construcción, 
  regularización y renovación, cubriendo el ciclo de vida completo de una obra. 
  El cliente puede personalizar su programa con una combinación de estos proyectos.
  \item \textbf{Hitos}: Un hito es un evento significativo dentro del proyecto. 
  Usualmente, lo asociamos a la entrega de recinto terminado, por ejemplo \emph{Cocina}.
  \item \textbf{Historias de usuario}: Esta es la lista de deseos de nuestro cliente. 
  Las historias de usuarios relacionadas se agrupan en el mismo Hito. Esto nos da la 
  oportunidad de pensar en la obra como un \emph{producto} enfocado en el cliente.
  \item \textbf{Sub-Tareas}: Una sub-tarea es la acción concreta que nuestros 
  especialista debe realizar para hacer realidad los deseos del cliente. Comúnmente, 
  son necesarias varias sub-tareas para completar una Historia de usuario.
\end{itemize}

\subsection{¿Cómo dividimos nuestro trabajo?}
Nosotros dividimos el trabajo en \emph{iteraciones} de \textbf{dos semanas}. Cada iteración 
tiene una meta, que llamamos \textbf{Meta quincenal}. Esta meta es establecida 
tanto por el cliente como por el equipo de trabajo.

\paragraph{}
Seleccionamos las historias de usuario relacionadas a la meta quincenal que se 
puedan completar en la iteración.

\paragraph{}
A veces hay historias de usuario que no se pueden incluír en la iteración, debido 
a circunstancias inciertas que solo se resolverán durante la iteración. Éstas 
historias son agrupadas en una lista llamada \textbf{Meta estrecha}.

\paragraph{}
Todo las demás historias son agrupadas en una lista llamada 
\textbf{Tareas pendientes}.